\documentclass[10pt,a4paper]{amsart}
\usepackage[margin=19mm]{geometry}
\usepackage{amsmath,amssymb,mathtools}
\usepackage[scr=rsfs,cal=euler]{mathalfa}
\usepackage{xfrac}
\usepackage[unicode,hidelinks,bookmarks=false]{hyperref}
\newcommand{\ie}{i.e.\ }
\newcommand{\cf}{cf.\ }
\newcommand{\resp}{respectively\ }
\newcommand{\Subsection}{\subsection}
\providecommand{\nobreakdash}{\nobreak\hbox{-}}
\setlength{\emergencystretch}{2em}
\multlinegap0pt
\usepackage{bm}
\usepackage{bbm}
\usepackage{dsfont}
\usepackage{xspace}
\usepackage{cancel}
\usepackage{mathtools}
\usepackage{amsthm}
\makeatletter
\@ifundefined{theorem}{\newtheorem{theorem}{Theorem}}{}
\makeatother
\newcommand{\op}{\oplus}
\newcommand{\perpdef}{\mathrel{\perp}}
\title[The first fatal axiom for weakened sequential products on fi]{The first fatal axiom for weakened sequential products on finite MV-effect algebras: Local obstruction, exact low-rank classification, and the rank-one boundary case}
\author{Joaquim Reizi Higuchi}
\date{Source: arXiv:2604.03324v1 (CC BY 4.0)}
\begin{document}
\maketitle

\noindent\emph{Source and licence.} The first fatal axiom for weakened sequential products on finite MV-effect algebras: Local obstruction, exact low-rank classification, and the rank-one boundary case. Version arXiv:2604.03324v1 (CC BY 4.0), distributed under the Creative Commons Attribution 4.0 International licence, as recorded in the arXiv licence metadata for this exact version and shown on the abstract page https://arxiv.org/abs/2604.03324v1. Full rights notice: licenses/arxiv-2604.03324-CC-BY-4.0.txt.\par\smallskip

\noindent\textit{Editorial scope.} Counted unit: the Theorem whose source label is \texttt{thm:B2-S123-classification} in the article (source0.tex lines 497--520, source label \texttt{thm:B2-S123-classification}), delivered together with the article's own complete original proof of it (source0.tex lines 522--583, one \texttt{proof} environment of 62 lines). No mathematics is written, repaired or re-derived: statement and proof are the article's own text line for line. Required context retained uncounted: none. Every definition, display and label of those lines is retained unchanged. Omitted from this package and not redistributed: 1--496, 521--521, 584--702. Nothing inside a retained excerpt is omitted or altered: every retained body is delivered exactly as the article's lines are written, so no omission receipt of any kind accompanies this package. Citations: the retained excerpts carry no citation command at all, so no bibliography entry of the article is transcribed and no reference is redistributed. Reference adaptation: none was needed and none was made; every reference inside the delivered unit resolves inside it, so no unresolved reference or citation is produced. Printed slips kept literal: no printed word, symbol, hypothesis or number of the counted statement or of its proof was altered. Layout and class: the article's own class is \texttt{article} with packages including geometry, fontenc, inputenc, lmodern, microtype, amsmath, amssymb, amsthm, mathtools, booktabs, enumitem, hyperref; the wrapper is the lane's pinned \texttt{amsart} editorial wrapper at 10pt on A4, which restates the theorem-like environments the retained text needs and adds the article's own macro definitions that the retained text needs (\texttt{\textbackslash op}, \texttt{\textbackslash perpdef}), with extra wrapper packages (bm, bbm, dsfont, xspace, cancel, mathtools). The delivered numbering is the unit's own. Assets: no retained span contains a figure environment, a graphics-inclusion command, a PDF-page inclusion, a table or a TikZ picture; no image file is copied into this package, the package declares no asset dependency and no image file is redistributed with it. Licence: version arXiv:2604.03324v1 of this article is distributed under the Creative Commons Attribution 4.0 International licence, as recorded in the arXiv licence metadata for this exact version and shown on the abstract page https://arxiv.org/abs/2604.03324v1. A full rights notice is delivered at licenses/arxiv-2604.03324-CC-BY-4.0.txt. Characters: every retained excerpt is pure ASCII, so the delivered engine, XeLaTeX with its default Latin Modern text font, reports no missing character.
\begin{theorem}[Exact \textup{(S1)}--\textup{(S3)} classification on $B_2$]\label{thm:B2-S123-classification}
A binary operation $\circ$ on $B_2$ satisfies \textup{(S1)}, \textup{(S2)}, and \textup{(S3)} if and only if there exist nonzero additive maps
\[
A,B\colon B_2\to B_2
\]
such that
\begin{equation}\label{eq:cross-zero-condition}
A(q)=0 \quad\Longleftrightarrow\quad B(p)=0
\end{equation}
and
\begin{equation}\label{eq:B2-row-form}
0\circ x=0,\qquad p\circ x=A(x),\qquad q\circ x=B(x),\qquad 1\circ x=x
\qquad (x\in B_2).
\end{equation}
Equivalently, if we write
\[
u:=A(p),\ v:=A(q),\ s:=B(p),\ t:=B(q),
\]
then $u\perpdef v$, $s\perpdef t$, $(u,v)\neq(0,0)$, $(s,t)\neq(0,0)$, and
\[
v=0 \quad\Longleftrightarrow\quad s=0.
\]
There are exactly $34$ such operations.
\end{theorem}

\begin{proof}
Assume first that $\circ$ satisfies \textup{(S1)}--\textup{(S3)}. Since $1\circ 0=0$, axiom \textup{(S3)} gives $0\circ 1=0$. Now
\[
1=x\op x'
\]
for every $x\in B_2$, so additivity of the left translation by $0$ yields
\[
0=0\circ 1=(0\circ x)\op(0\circ x').
\]
By positivity, $0\circ x=0$ for all $x\in B_2$. Thus the row at $0$ is the zero map.

Let
\[
A(x):=p\circ x,
\qquad
B(x):=q\circ x.
\]
By \textup{(S1)}, both $A$ and $B$ are additive. By \textup{(S2)}, the row at $1$ is the identity map.

We claim that $A$ and $B$ are nonzero. If $A=0$, then in particular $p\circ 1=0$, so \textup{(S3)} gives
\[
1\circ p=p\circ 1=0,
\]
contradicting \textup{(S2)}. Hence $A\neq 0$. The same argument shows $B\neq 0$.

Finally,
\[
A(q)=p\circ q=0
\]
if and only if
\[
q\circ p=0,
\]
and the latter equality is exactly $B(p)=0$. Therefore \eqref{eq:cross-zero-condition} holds, and \eqref{eq:B2-row-form} follows from the definitions.

Conversely, assume that $A$ and $B$ are nonzero additive maps satisfying \eqref{eq:cross-zero-condition}, and define $\circ$ by \eqref{eq:B2-row-form}. Then \textup{(S1)} holds because each left translation is additive, and \textup{(S2)} holds because the row at $1$ is the identity.

It remains to prove \textup{(S3)}. Suppose that $a\circ b=0$.

If $a=0$, then $b\circ a=b\circ 0=0$ because every additive map sends $0$ to $0$. If $a=1$, then $b=1\circ b=a\circ b=0$, so again $b\circ a=0$. Thus only $a\in\{p,q\}$ needs attention.

If $b=0$, then $b\circ a=0$ as above. If $b=1$, then $a\circ 1$ is either $A(1)$ or $B(1)$, and this is nonzero because $A$ and $B$ are nonzero additive maps on $B_2$. Hence the case $b=1$ cannot occur under the assumption $a\circ b=0$.

If $(a,b)=(p,p)$ or $(q,q)$, then $b\circ a=a\circ a=0$ trivially. If $(a,b)=(p,q)$, then $A(q)=0$, so \eqref{eq:cross-zero-condition} gives $B(p)=0$, hence
\[
b\circ a=q\circ p=B(p)=0.
\]
The case $(a,b)=(q,p)$ is symmetric. Therefore $a\circ b=0$ always implies $b\circ a=0$, so \textup{(S3)} holds.

For the counting statement, note that a nonzero additive map $A$ is determined by an orthogonal pair $(u,v)=(A(p),A(q))\neq(0,0)$. There are exactly three such pairs with $v=0$, namely
\[
(p,0),\ (q,0),\ (1,0),
\]
and exactly five such pairs with $v\neq 0$, namely
\[
(0,p),\ (0,q),\ (0,1),\ (p,q),\ (q,p).
\]
Likewise, a nonzero additive map $B$ is determined by an orthogonal pair $(s,t)=(B(p),B(q))\neq(0,0)$; there are three possibilities with $s=0$ and five with $s\neq 0$. Condition \eqref{eq:cross-zero-condition} matches the two zero-patterns, so the total number of operations is
\[
3\cdot 3+5\cdot 5=34.
\]
\end{proof}

\end{document}
